% paralleles.tex — f166r, vv. 294–336 : loci paralleli. Généré par .travail/f166r_lp_redac2/rediger.py
% à partir des collectes vérifiées (paralleles_294-315.json, paralleles_316-336.json) et des parallèles
% d'Austin–Olson 2004 retrouvés dans les sources ; règles : edition/CONVENTIONS.md § 4
\lp{294 δούλοις γὰρ οὐκ ἔξεστ᾽ ἀκούειν]}{Ar.~\textit{Thesm.}~330 εὐγενεῖς γυναῖκες ; Ar.~\textit{Thesm.}~541 εἰ γὰρ οὔσης / παρρησίας κἀξὸν λέγειν ὅσαι πάρεσμεν ἀσταί ; Trag. adesp. ap. Ph.~\textit{Prob.}~48 «οὐ γὰρ μετεῖναι τῶν νόμων δούλοις ἔφυ.» καὶ πάλιν· «δοῦλος πέφυκας, οὐ μέτεστί σοι λόγου.» ; \textit{idem} M.Ant.~11.30 (om. οὐ γὰρ μετεῖναι … καὶ πάλιν) ; Is.~6.49–50 ἡ δὲ τούτων μήτηρ, οὕτως ὁμολογουμένως οὖσα δούλη … ἣν οὔτε παρελθεῖν εἴσω τοῦ ἱεροῦ ἔδει οὔτ᾽ ἰδεῖν τῶν ἔνδον οὐδέν, οὔσης τῆς θυσίας ταύταις ταῖς θεαῖς ἐτόλμησε συμπέμψαι τὴν πομπὴν καὶ εἰσελθεῖν εἰς τὸ ἱερὸν καὶ ἰδεῖν ἃ οὐκ ἐξῆν αὐτῇ ; D.~59.85 οὐκ ἔξεστιν αὐτῇ ἐλθεῖν εἰς οὐδὲν τῶν ἱερῶν τῶν δημοτελῶν, εἰς ἃ καὶ τὴν ξένην καὶ τὴν δούλην [ἐλθεῖν] ἐξουσίαν ἔδοσαν οἱ νόμοι καὶ θεασομένην καὶ ἱκετεύσουσαν εἰσιέναι ; Plu.~\textit{Arist.}~21.2 δούλῳ γὰρ οὐδενὸς ἔξεστι τῶν περὶ τὴν διακονίαν ἐκείνην προσάψασθαι ; Plu.~\textit{Quaest.\,Rom.}~16 (267d–e) διὸ καὶ παρ᾽ ἡμῖν ἐν Χαιρωνείᾳ πρὸ τοῦ σηκοῦ τῆς Λευκοθέας ὁ νεωκόρος λαβὼν μάστιγα κηρύσσει μὴ δοῦλον εἰσιέναι μὴ δούλαν ; Ath.~6.262c φησὶ γὰρ Μακαρεὺς ἐν τῇ τρίτῃ Κῳακῶν ὅτι ὁπόταν τῇ Ἥρᾳ θύωσιν οἱ Κῷοι οὔτε εἴσεισιν εἰς τὸ ἱερὸν δοῦλος οὔτε γεύεταί τινος τῶν παρεσκευασμένων ; Zen.~4.33 Θύραζε Κᾶρες, οὐκ ἔτ᾽ Ἀνθεστήρια: οἱ μὲν διὰ πλῆθος οἰκετῶν Καρικῶν εἰρῆσθαι φασὶν, ὡς ἐν τοῖς Ἀνθεστηρίοις εὐωχουμένων αὐτῶν καὶ οὐκ ἐργαζομένων ; \textit{idem} Diogenian.~5.23 (om. οἱ μέν et Καρικῶν εἰρῆσθαι φασίν ; εὐωχουμένων καὶ αὐτῶν pro εὐωχουμένων αὐτῶν καὶ οὐκ ἐργαζομένων), Erasm.~\textit{Adag.}~1.2.65 (n°~165) (Foras Cares, non amplius anthisteria … Cares autem servos appellabant) ; Ar.~\textit{Ach.}~248–250 τήνδε τὴν πομπὴν ἐμὲ / πέμψαντα καὶ θύσαντα μετὰ τῶν οἰκετῶν / ἀγαγεῖν τυχηρῶς τὰ κατ᾽ ἀγροὺς Διονύσια ; Ar.~\textit{Thesm.}~537 αὐταί τε καὶ τὰ δουλάρια τέφραν ποθὲν λαβοῦσαι ; Ar.~\textit{Thesm.}~608–609 ἡδὶ δὲ δὴ τίς ἐστιν ἡ τὸ παιδίον / ἔχουσα; — τίτθη νὴ Δί᾽ ἐμή ; Ar.~\textit{Thesm.}~728 ἴωμεν ἐπὶ τὰς κληματίδας ὦ Μανία ; Is.~8.19 αἵ τε γυναῖκες αἱ τῶν δημοτῶν μετὰ ταῦτα προὔκριναν αὐτὴν μετὰ τῆς Διοκλέους γυναικὸς τοῦ Πιθέως ἄρχειν εἰς τὰ Θεσμοφόρια ; Aeschin.~1.138 ταῦτα τοῖς δούλοις ἀπεῖπον μὴ ποιεῖν. δοῦλον, φησὶν ὁ νόμος, μὴ γυμνάζεσθαι μηδὲ ξηραλοιφεῖν ἐν ταῖς παλαίστραις ; Ath.~14.639d Κῷοι δὲ τοὐναντίον δρῶσιν, ὡς ἱστορεῖ Μακαρεὺς ἐν τρίτῳ Κῳακῶν· ὅταν γὰρ τῇ Ἥρᾳ θύωσιν, δοῦλοι οὐ παραγίνονται ἐπὶ τὴν εὐωχίαν ; Eur.~\textit{Hec.}~234–235 εἰ δ᾽ ἔστι τοῖς δούλοισι τοὺς ἐλευθέρους / μὴ λυπρὰ μηδὲ καρδίας δηκτήρια / ἐξιστορῆσαι ; Eur.~\textit{Hel.}~733 ἄλλων τ᾽ ἀκούειν δοῦλον ὄντα τῶν πέλας}
\lp{294 οὐκ ἔξεστ᾽ ἀκούειν (βέβηλοι)]}{Ar.~\textit{Ran.}~354–355 εὐφημεῖν χρὴ κἀξίστασθαι τοῖς ἡμετέροισι χοροῖσιν, / ὅστις ἄπειρος τοιῶνδε λόγων ἢ γνώμῃ μὴ καθαρεύει ; Pl.~\textit{Smp.}~218b οἱ δὲ οἰκέται, καὶ εἴ τις ἄλλος ἐστὶν βέβηλός τε καὶ ἄγροικος, πύλας πάνυ μεγάλας τοῖς ὠσὶν ἐπίθεσθε ; Isoc.~4.157–158 Εὐμολπίδαι δὲ καὶ Κήρυκες ἐν τῇ τελετῇ τῶν μυστηρίων διὰ τὸ τούτων μῖσος καὶ τοῖς ἄλλοις βαρβάροις εἴργεσθαι τῶν ἱερῶν, ὥσπερ τοῖς ἀνδροφόνοις, προαγορεύουσιν ; Ar.~\textit{Pax}~966–969 οὐχ αἱ γυναῖκές γ᾽ ἔλαβον. — ἀλλ᾽ εἰς ἑσπέραν / δώσουσιν αὐταῖς ἅνδρες. — ἀλλ᾽ εὐχώμεθα. / τίς τῇδε; ποῦ ποτ᾽ εἰσὶ πολλοὶ κἀγαθοί ; Ar.~\textit{Ran.}~369–370 τούτοις αὐδῶ καὖθις ἀπαυδῶ καὖθις τὸ τρίτον μάλ᾽ ἀπαυδῶ / ἐξίστασθαι μύσταισι χοροῖς ; Eur.~fr.~648~N.² (\textit{Protesilaus}) οὐ γὰρ θέμις βέβηλον ἅπτεσθαι δόμων ; Call.~\textit{Ap.}~2 ἑκὰς ἑκὰς ὅστις ἀλιτρός ; Call.~\textit{Cer.}~3–5 τὸν κάλαθον κατιόντα χαμαὶ θασεῖσθε βέβαλοι, / μηδ᾽ ἀπὸ τῶ τέγεος μηδ᾽ ὑψόθεν αὐγάσσησθε / μὴ παῖς μηδὲ γυνὰ μηδ᾽ ἃ κατεχεύατο χαίταν ; Verg.~\textit{Aen.}~6.258–259 Procul O procul este, profani, / conclamat vates ; \textit{idem} Tosi n°~224 (om. conclamat vates) ; Aristid.~\textit{Or.}~46 (p.~129 Jebb) λέγω δὲ ἤδη ταῦτα πρὸς οὓς ἔξεστι. τῶν δὲ βεβήλων ὀλίγος ὁ λόγος. πάντως οὐδὲν δεῖ πύλας αὐτοὺς ἐπιθέσθαι τοῖς ὠσὶν, ἀλλ᾽ ἐπίκεινται πάλαι ; App. Prov.~4.90 Τίς τῇδε: οἱ σπένδοντες ἔλεγον τοῦτο ἀντὶ τοῦ, τίς πάρεστι. οἱ δὲ παρόντες ἐπευφημούμεναι ἀπεκρίνοντο, Καλοὶ κἀγαθοί. τοῦτο δὲ ἐποίουν [οἱ σπένδοντες], ἵνα οἱ συνειδότες ἑαυτοῖς τι ἄτοπον, ἐκχωροῖεν τῶν σπονδῶν ; Erasm.~\textit{Adag.}~3.5.18 (n°~2418) Ostium obdite profani. Θύρας ἐπίθεσθε βέβηλοι. i. e. Fores occludite quique profani. … Jubet ut profani ostium auribus imponant, ne quid arcani sermonis possit ingredi ; Hor.~\textit{Carm.}~3.1.1 (Tosi n°~223) Odi profanum volgus et arceo}
\lp{294 οὐκ ἔξεστ᾽ ἀκούειν (Θεσμοφόρια)]}{Ar.~\textit{Thesm.}~472 αὐταὶ γάρ ἐσμεν, κοὐδεμί᾽ ἔκφορος λόγου ; Ar.~\textit{Thesm.}~1150–1151 οὗ δὴ ἀνδράσιν οὐ θεμιτεύετον εἰσορᾶν / ὄργια σεμνὰ θεοῖν ; Hdt.~2.171.2 καὶ τῆς Δήμητρος τελετῆς πέρι, τὴν οἱ Ἕλληνες θεσμοφόρια καλέουσι, καὶ ταύτης μοι πέρι εὔστομα κείσθω, πλὴν ὅσον αὐτῆς ὁσίη ἐστὶ λέγειν ; Hsch.~δ~2036 δίωγμα· θυσία τις Ἀθήνησιν ἐν ἀπορρήτῳ τελούμενη ὑπὸ τῶν γυναικῶν ἐν τοῖς Θεσμοφορίοις. τὸ αὐτὸ καὶ ἀποδίωγμα ὕστερον ἐκλήθη ; Suid.~θ~272 Θεσμοφόρος. ὅτι Βάττος ὁ Κυρήνην κτίσας τῆς Θεσμοφόρου τὰ μυστήρια ἐγλίχετο μαθεῖν καὶ προσῆγε βίαν λίχνοις ὀφθαλμοῖς χαριζόμενος ; Hdt.~6.134–135 ἕρκος θεσμοφόρου Δήμητρος ὑπερθορεῖν … τὰ ἐς ἔρσενα γόνον ἄρρητα ἱρὰ ἐκφήνασαν Μιλτιάδῃ}
\lp{295–311 εὐφημία ἔστω … χαίρωμεν]}{Ar.~\textit{Av.}~864–888 εὔχεσθε τῇ Ἑστίᾳ τῇ ὀρνιθείῳ καὶ τῷ ἰκτίνῳ τῷ ἑστιούχῳ καὶ ὄρνισιν Ὀλυμπίοις καὶ Ὀλυμπίῃσι πᾶσι καὶ πάσῃσιν ; Aeschin.~1.23 ἐπειδὰν τὸ καθάρσιον περιενεχθῇ καὶ ὁ κῆρυξ τὰς πατρίους εὐχὰς εὔξηται … καὶ μετὰ ταῦτα ἐπερωτᾷ ὁ κῆρυξ· τίς ἀγορεύειν βούλεται τῶν ὑπὲρ πεντήκοντα ἔτη γεγονότων ; Din.~2.14 καὶ ὁ μὲν νόμος εὐξάμενον κελεύει τὸν κήρυκα μετ᾽ εὐφημίας πολλῆς, οὕτως ὑμῖν τὸ βουλεύεσθαι περὶ τῶν πραγμάτων παραδιδόναι ; Ar.~\textit{Av.}~1035 ἐὰν δ᾽ ὁ Νεφελοκοκκυγιεὺς τὸν Ἀθηναῖον ἀδικῇ — ; Ar.~\textit{Av.}~1661 νόθῳ δὲ μὴ εἶναι ἀγχιστείαν παίδων ὄντων ; Archipp.~fr.~27~K.-A. ἀποδοῦναι δ᾽ ὅσα ἔχομεν ἀλλήλων, ἡμᾶς μὲν τὰς Θρᾴττας καὶ Ἀθερίνην τὴν αὐλητρίδα καὶ Σηπίαν ; Ar.~\textit{Eq.}~941 εὖ γε νὴ τὸν Δία καὶ τὸν Ἀπόλλω καὶ τὴν Δήμητρα ; Ar.~\textit{Av.}~1040–1041 χρῆσθαι Νεφελοκοκκυγιᾶς τοῖσδε τοῖς μέτροισι καὶ / σταθμοῖσι καὶ ψηφίσμασι καθάπερ Ὀλοφύξιοι ; Ar.~\textit{Av.}~1046 καλοῦμαι Πισθέταιρον ὕβρεως ἐς τὸν Μουνιχιῶνα μῆνα ; D.~54.39 τοὺς ὄρχεις τοὺς ἐκ τῶν χοίρων, οἷς καθαίρουσιν ὅταν εἰσιέναι μέλλωσιν}
\lp{295–296 εὐφημία ἔστω]}{\textit{Test.} Suid.~ε~3797 Εὐφημία ἔστω: εὐφημίᾳ χρῶνται οἱ κωμικοί, ἐπειδὰν εὐχὴν ἢ ψήφισμα εἰσάγωσιν ; \textit{idem} Σ Ar.~\textit{Thesm.}~295 (Dübner, p.~266,~25–26) (πεζῇ εὐφημίᾳ ; om. Εὐφημία ἔστω) ; Ar.~\textit{Ach.}~237–238 εὐφημεῖτε, εὐφημεῖτε. — σῖγα πᾶς. ἠκούσατ᾽ ἄνδρες ἆρα τῆς εὐφημίας ; Ar.~\textit{Eq.}~1316–1317 εὐφημεῖν χρὴ καὶ στόμα κλῄειν καὶ μαρτυριῶν ἀπέχεσθαι, / καὶ τὰ δικαστήρια συγκλῄειν οἷς ἡ πόλις ἥδε γέγηθεν ; Ar.~\textit{Nub.}~263 εὐφημεῖν χρὴ τὸν πρεσβύτην καὶ τῆς εὐχῆς ἐπακούειν ; Ar.~\textit{Vesp.}~868 εὐφημία μὲν πρῶτα νῦν ὑπαρχέτω ; Ar.~\textit{Pax}~433–435 σπονδὴ σπονδή· / εὐφημεῖτε εὐφημεῖτε. / σπένδοντες εὐχόμεσθα τὴν νῦν ἡμέραν ; Ar.~\textit{Av.}~959 εὐφημία ᾽στω ; Ar.~\textit{Thesm.}~39–40 εὔφημος πᾶς ἔστω λαός, / στόμα συγκλῄσας ; Ar.~\textit{Ran.}~354 εὐφημεῖν χρὴ κἀξίστασθαι τοῖς ἡμετέροισι χοροῖσιν ; Eur.~\textit{Ba.}~68–70 τίς ὁδῷ τίς ὁδῷ; τίς; / μελάθροις ἔκτοπος ἔστω, στόμα τ᾽ εὔφη-/μον ἅπας ἐξοσιούσθω ; Eur.~\textit{Hec.}~530–533 σιγὴν Ἀχαιῶν παντὶ κηρῦξαι στρατῷ. / κἀγὼ καταστὰς εἶπον ἐν μέσοις τάδε· / Σιγᾶτ᾽, Ἀχαιοί, σῖγα πᾶς ἔστω λεώς, / σίγα σιώπα ; Eur.~\textit{IA}~1563–1564 στὰς δ᾽ ἐν μέσῳ Ταλθύβιος, ᾧ τόδ᾽ ἦν μέλον, / εὐφημίαν ἀνεῖπε καὶ σιγὴν στρατῷ ; \textit{Il.}~9.171–172 φέρτε δὲ χερσὶν ὕδωρ, εὐφημῆσαί τε κέλεσθε, / ὄφρα Διὶ Κρονίδῃ ἀρησόμεθ᾽ ; Hor.~\textit{Carm.}~3.1.2 favete linguis ; \textit{idem} Tosi n°~2083 (add. εὐφημεῖτε) ; Th.~6.32.1 τῇ μὲν σάλπιγγι σιωπὴ ὑπεσημάνθη, εὐχὰς δὲ τὰς νομιζομένας πρὸ τῆς ἀναγωγῆς οὐ κατὰ ναῦν ἑκάστην, ξύμπαντες δὲ ὑπὸ κήρυκος ἐποιοῦντο ; Pl.~\textit{Lg.}~800e–801a νικᾷ γὰρ πάσαισι ταῖς ψήφοις οὗτος ὁ νόμος. τίς δὴ μετ᾽ εὐφημίαν δεύτερος ἂν εἴη νόμος μουσικῆς; ἆρ᾽ οὐκ εὐχὰς εἶναι τοῖς θεοῖς οἷς θύομεν ἑκάστοτε ; Σ Ar.~\textit{Thesm.}~295 (Dübner, p.~266,~24) γυνὴ μιμουμένη κήρυκα ; Hsch.~ε~7264 εὐφημίαν· σιωπήν (Soph. Tr. 178). εὐλογίαν (Aeschin. 1, 169) ; Suid.~ε~3792 Εὐφήμει: καλὰ λάλει, ἢ ἐπαίνει. Ἀριστοφάνης· εὔφημος πᾶς ἔστω λεὼς στόμα συγκλείσας. ἐνδημεῖ γὰρ θίασος Μουσῶν ἔνδον μελάθρων. Εὐφήμει οὖν σιώπα, μὴ ἀκαιρολόγει ; \textit{idem} Hsch.~ε~7269 (εὐφημεῖ· καλὰ λαλεῖ, ἢ ἐπαινεῖ tantum), Phot.~ε~2360 (εὐφημεῖ· καλὰ λαλεῖ, ἐπαινεῖ tantum), Phot.~ε~2361 (εὐφήμει· σιώπα, μὴ ἀκαιρολόγει tantum) ; Suid.~ε~3793 εὔφημον οὖν τὴν σιωπηλήν· τὸ γὰρ εὐφημεῖν ἐπὶ τοῦ σιωπᾶν τίθεται ; Ar.~\textit{Ach.}~123 σίγα, κάθιζε ; Ar.~\textit{Nub.}~297 ἀλλ᾽ εὐφήμει· μέγα γάρ τι θεῶν κινεῖται σμῆνος ἀοιδαῖς ; Ar.~\textit{Pax}~96–97 εὐφημεῖν χρὴ καὶ μὴ φλαῦρον / μηδὲν γρύζειν ἀλλ᾽ ὀλολύζειν ; Ar.~\textit{Pax}~1316 εὐφημεῖν χρὴ καὶ τὴν νύμφην ἔξω τινὰ δεῦρο κομίζειν ; Ar.~\textit{Thesm.}~381–382 σίγα σιώπα, πρόσεχε τὸν νοῦν· χρέμπτεται γὰρ ἤδη / ὅπερ ποιοῦσ᾽ οἱ ῥήτορες ; Ar.~\textit{Ran.}~1274 εὐφαμεῖτε· μελισσονόμοι δόμον Ἀρτέμιδος πέλας οἴγειν ; Eup.~fr.~413~K.-A. (ap. Phot.~ε~2149) εὖ ἔχειν στόμα· τὸ εὐφημεῖν. οὕτως Εὔπολις ; Eur.~\textit{IA}~1467–1469 ὑμεῖς δ᾽ ἐπευφημήσατ᾽, ὦ νεάνιδες, / παιᾶνα τἠμῇ συμφορᾷ Διὸς κόρην / Ἄρτεμιν· ἴτω δὲ Δαναΐδαις εὐφημία ; Eur.~\textit{IT}~123–125 εὐφαμεῖτ᾽, ὦ / πόντου δισσὰς συγχωρούσας / πέτρας Ἀξείνου ναίοντες ; Eur.~\textit{Ion}~98 στόμα τ᾽ εὔφημον φρουρεῖν ἀγαθόν ; Eur.~\textit{Supp.}~668–669 κῆρυξ δὲ Θησέως εἶπεν ἐς πάντας τάδε· / Σιγᾶτε, λαοί· σῖγα, Καδμείων στίχες ; Eur.~fr.~1087~N.² εὐφημία γὰρ εὐκολώτατος πόνων ; \textit{idem} Suid.~ε~3795 (praem. εὐφημίας χάριν προτιμᾷ τὴν Λακεδαίμονα τὸ θεῖον ἤπερ τῶν Ἀθηναίων βωμοὺς σύμπαντας … ; add. … καὶ Εὐριπίδης· εὐφημία γὰρ παρὰ σπονδαῖσι κάλλιστον) ; Aesch.~\textit{Eum.}~1035 εὐφαμεῖτε δέ, χωρῖται ; Soph.~\textit{Trach.}~178 εὐφημίαν νῦν ἴσχ᾽ ; Call.~\textit{Ap.}~17 εὐφημεῖτ᾽ ἀίοντες ἐπ᾽ Ἀπόλλωνος ἀοιδῆι ; Verg.~\textit{Aen.}~5.71 Ore favete omnes ; [Pl.]~\textit{Alc.2}~149b φησὶν ἂν βούλεσθαι αὑτῷ τὴν Λακεδαιμονίων εὐφημίαν εἶναι μᾶλλον ἢ τὰ σύμπαντα τῶν Ἑλλήνων ἱερά. … τήν γ᾽ οὖν εὐφημίαν οὐκ ἄλλην τινά μοι δοκεῖ λέγειν ὁ θεὸς ἢ τὴν εὐχὴν αὐτῶν ; Plu.~\textit{Cor.}~25.2 ὅταν γὰρ ἄρχοντες ἢ ἱερεῖς πράττωσί τι τῶν θείων, ὁ κῆρυξ πρόεισι μεγάλῃ φωνῇ βοῶν, Ὃκ ἄγε ; Ath.~12.538a (Ephipp. FGrHist~126 F~5) ἐθυμιᾶτο δὲ αὐτῷ σμύρνα καὶ τὰ ἄλλα θυμιάματα· εὐφημία τε καὶ σιγὴ κατεῖχε πάντας ὑπὸ δέους τοὺς παρόντας ; Σ Il.~9.171 εὐφημῆσαι] εὔφημον καὶ εὔτακτον φθέγξασθαι· οὐδὲ γὰρ ἐπὶ τοῦ ὑμνῆσαι εὕρηται παρὰ τῷ Ὁμήρῳ ; Hsch.~ε~7270 εὐφημῆσαι· βοῆσαι μετ᾽ εὐφημίας (Ι 171) ; Hsch.~ι~132 ἰβυλῆνας· τοὺς εὐφημοῦντας· ἴβυς γὰρ ἡ εὐφημία ; Phot.~ε~2339 εὐστομεῖν· τὸ εὐφημεῖν· καὶ οἱ Ἴωνες «εὔστομα κείσθω» φασίν ; Phot.~ε~2362 εὐφημίαν· ἐπὶ τῆς σιγῆς ἔλεγον ; Apostol.~8.84a (CPG II, p.~455) Θειότερον εὐφημία κακηγορίας ; Erasm.~\textit{Adag.}~3.5.18 (n°~2418) Ante sacrificium praeco clamore solebat jubere populum favere linguis, h. e. verba boni ominis loqui … Hoc verbum non ut plerique existimant à favore trahitur, sed imperatur silentium, ut rite peragi possit sacrum, nulla voce mala obstrepente ; Aesch.~\textit{Ag.}~1247 εὔφημον, ὦ τάλαινα, κοίμησον στόμα ; Pl.~\textit{Phd.}~117e ἐγὼ μέντοι οὐχ ἥκιστα τούτου ἕνεκα τὰς γυναῖκας ἀπέπεμψα … καὶ γὰρ ἀκήκοα ὅτι ἐν εὐφημίᾳ χρὴ τελευτᾶν}
\lp{295–296 εὐφημία ἔστω, εὐφημία ἔστω]}{Eur.~\textit{Hipp.}~58,~61 ἕπεσθ᾽ ἀείδοντες ἕπεσθε … πότνια πότνια σεμνοτάτα ; Verg.~\textit{Aen.}~6.46 deus, ecce, deus!}
\lp{297–301 εὔχεσθε ταῖν Θεσμοφόροιν … καὶ ταῖς Χάρισιν]}{Hicks–Hill~119.5–11 (= IG II²~112) εὔξασθαι μὲν τὸγ κήρυκα αὐτίκα μάλα τῶι Διὶ τῶι Ὀλυμπίωι καὶ τῆι Ἀθηνᾶι τῆι Πολιάδι καὶ τῆι Δήμητρι καὶ τῆι Κόρηι καὶ τοῖς δώδεκα [θ]εοῖς καὶ ταῖς σεμναῖς θεαῖς}
\lp{297 εὔχεσθε]}{Ar.~\textit{Thesm.}~331 εὔχεσθε τοῖς θεοῖσι τοῖς Ὀλυμπίοις ; Ar.~\textit{Thesm.}~350–351 ταῖς δ᾽ ἄλλαισιν ὑμῖν τοὺς θεοὺς / εὔχεσθε πάσαις πολλὰ δοῦναι κἀγαθά ; Hes.~\textit{Op.}~465–466 εὔχεσθαι δὲ Διὶ χθονίῳ Δημήτερί θ᾽ ἁγνῇ, / ἐκτελέα βρίθειν Δημήτερος ἱερὸν ἀκτήν ; Ar.~\textit{Ran.}~885 εὔχεσθε δὴ καὶ σφώ τι πρὶν τἄπη λέγειν ; \textit{Il.}~7.194 τόφρ᾽ ὑμεῖς εὔχεσθε Διὶ Κρονίωνι ἄνακτι}
\lp{297 ταῖν Θεσμοφόροιν]}{Ar.~\textit{Thesm.}~282 ἀλλ᾽ ὦ περικαλλεῖ Θεσμοφόρω δέξασθέ με ; Ar.~\textit{Thesm.}~1156 Θεσμοφόρω πολυποτνία ; Ar.~\textit{Thesm.}~1229–1231 τὼ Θεσμοφόρω δ᾽ / ἡμῖν ἀγαθὴν / τούτων χάριν ἀνταποδοίτην ; Soph.~\textit{OC}~1050 οὗ πότνιαι σεμνὰ τιθηνοῦνται τέλη ; Σ Ar.~\textit{Thesm.}~83 (Dübner, p.~265,~12–14) κἀν θεσμοφόροιν: Ἀττικὴ σύνταξις καὶ δυϊκή, ἀντὶ τοῦ ἐν τῷ τῶν θεσμοφόρων ἱερῷ, Δήμητρος καὶ Περσεφόνης ; Σ Soph.~\textit{OC}~1050 οὗ πότνιαι σεμναί: αἱ Θεσμοφόροι ; Σ Luc.~\textit{DMeretr.}~2.1 (Rabe, p.~275–276) Θεσμοφορία δὲ καλεῖται, καθότι θεσμοφόρος ἡ Δημήτηρ κατονομάζεται τιθεῖσα νόμους ἤτοι θεσμούς, καθ᾽ οὓς τὴν τροφὴν πορίζεσθαί τε καὶ κατεργάζεσθαι ἀνθρώπους δέον ; Hsch.~θ~384 θεσμοφόρος· ὁ θεσμοὺς φέρων, ἤτοι νόμους ; Ar.~\textit{Av.}~1519 ἀλλ᾽ ὡσπερεὶ Θεσμοφορίοις νηστεύομεν ; Ar.~\textit{Thesm.}~83 κἀν Θεσμοφόροιν μέλλουσι περί μου τήμερον ; Ar.~\textit{Eccl.}~443 ἐκ Θεσμοφόροιν ἑκάστοτ᾽ αὐτὰς ἐκφέρειν ; Paus.~1.31.1 Ἀλιμουσίοις ⟨μὲν⟩ Θεσμοφόρου Δήμητρος καὶ Κόρης ἐστὶν ἱερόν ; Hsch.~α~5234 ἄνοδος· ἀνάβασις. ἡ ἑνδέκατη τοῦ Πυανεψιῶνος, ὅτε αἱ γυναῖκες ἀνέρχονται εἰς Θεσμοφόριον, οὕτω καλεῖται ; Hsch.~θ~383 θεσμοφόρια· τὰ μυστήρια τῆς Δήμητρος. ἢ ὄνομα ἑορτῆς ; Hsch.~κ~3098 κνέωρον· φυτόν τι, ὃ τοῖς Θεσμοφορίοις ὑποστόρνυται, καὶ ὃ εἰς κάθαρσιν χρῶνται. καὶ γυναικεῖον μόριον (com. ad fr. 369 K.–A.) ; Hsch.~τ~1440 τρίτη Θεσμοφορίων· ζητεῖται πῶς ἅμα μὲν λέγει τρίτην Θεσμοφορίων εἶναι, ἅμα δὲ μέσην, τεττάρων οὐσῶν ἡμερῶν (Ar. Thesm. 80) ; Synag.~θ~63 θεσμοφόρος· ὁ τοὺς νόμους ἐπιφερόμενος ; \textit{idem} Phot.~θ~135, Suid.~θ~271 (add. καὶ ἐν Ἐπιγράμμασι· ὡμολόγησ᾽ ἥξειν εἰς νύκτα μοι ἐπιβόητος, καὶ σεμνὴν ὤμοσε Θεσμοφόρον) ; Phot.~θ~134 Θεσμοφορίων ἡμέραι· δ΄· δεκάτη Θεσμοφόρια, ἑνδεκάτη κάθοδος, δωδεκάτη νηστεία, τρισκαιδεκάτη καλλιγένεια ; Hsch.~ζ~145 ζημία· θυσία τις ἀποδιδομένη ὑπὲρ τῶν γινομένων ἐν Θεσμοφορίοις ; Hsch.~τ~1374 τριήμερος· Θεσμοφόρια ὑπὸ Λακώνων}
\lp{297 ταῖν Θεσμοφόροιν (τὼ θεώ)]}{Ar.~\textit{Eccl.}~155–158 ἐμοὶ μὲν οὐ δοκεῖ μὰ τὼ θεώ. — μὰ τὼ θεώ; τάλαινα ποῦ τὸν νοῦν ἔχεις; … μὰ Δί᾽ ἀλλ᾽ ἀνὴρ ὢν τὼ θεὼ κατώμοσας ; Hsch.~ν~543 νὴ τὼ θεώ· ὅρκος κατὰ Δήμητρος καὶ Κόρης ; Phot.~μ~147 μὰ τὼ θεώ (Ar. Eccl. 155)· γυναικεῖος ὅρκος· δυϊκῶς δὲ ὀμνύουσι τὴν Κόρην καὶ τὴν Δήμητραν· ἀνδράσι δὲ οὐ πρέπει τοῦτον ὀμνύναι ; Ar.~\textit{Thesm.}~285 τὸ πόπανον, ὅπως λαβοῦσα θύσω τοῖν θεοῖν ; Ar.~\textit{Thesm.}~383 φιλοτιμίᾳ μὲν οὐδεμιᾷ μὰ τὼ θεὼ ; Σ Luc.~\textit{DMeretr.}~7.4 κατὰ ταῖν θεαῖν] Κόρη καὶ Δήμητηρ αἱ θεαί, ἡ δὲ Ἀθηνᾶ ἡ Πολιάς. ὅρκος δὲ τοῦτο Ἀθηναίων, ὡς Λακεδαιμονίων οἱ Διόσκουροι ; Luc.~\textit{Pisc.}~33 ἤν τινα καὶ τῶν μεμυημένων ἰδὼν ἐξαγορεύοντα ταῖν θεαῖν τὰ ἀπόρρητα ; Alciphr.~3.26 μυστήριον ἐν αὐταῖς [σ]τρέφεται 〈τῶν〉 ταῖν θεαῖν ταῖν Ἐλευσινίαιν ἀσφαλέστερον}
\lp{[298] τῇ Δήμητρι καὶ τῇ Κόρῃ]}{D.H.~\textit{Din.}~11 διαδικασία Ἀθμονεῦσι περὶ τῆς μυρρίνης καὶ τῆς μίλακος· εὔχομαι δὴ τῇ Δήμητρι καὶ τῇ Κόρῃ ; Ar.~\textit{Ran.}~337 ὦ πότνια πολυτίμητε Δήμητρος κόρη ; Hdt.~8.65.4 τὴν δὲ ὁρτὴν ταύτην ἄγουσι Ἀθηναῖοι ἀνὰ πάντα ἔτεα τῇ Μητρὶ καὶ τῇ Κούρῃ ; Σ Ar.~\textit{Thesm.}~1230 (Dübner, p.~272) ἡ Δημήτηρ καὶ ἡ Περσεφόνη ; Eur.~\textit{Supp.}~290 σεμναῖσι Δηοῦς ἐσχάραις παρημένη ; Zen.~2.22 Ἀνθ᾽ Ἑρμιόνος; εἴρηται ἡ παροιμία ἐπὶ τῶν σωζόντων τοὺς ἱκέτας· ἐπειδὴ ἐν Ἑρμιόνῃ τῆς Πελοποννήσου ἱερὸν ἦν Κόρης καὶ Δήμητρος, ἀσφάλειαν παρέχον τοῖς καταφεύγουσιν}
\lp{299 καὶ τῷ Πλούτῳ]}{Ar.~\textit{Plut.}~727 μετὰ τοῦτο τῷ Πλούτωνι παρεκαθέζετο ; Ar.~fr.~488~Kock (\textit{Tagenistae}) καὶ μὴν πόθεν Πλούτων γ᾽ ἂν ὠνομάζετο, / εἰ μὴ τὰ βέλτιστ᾽ ἔλαχεν ; Hes.~\textit{Th.}~969–974 Δημήτηρ μὲν Πλοῦτον ἐγείνατο, δῖα θεάων, / Ἰασίων᾽ ἥρωι μιγεῖσ᾽ ἐρατῇ φιλότητι / νειῷ ἔνι τριπόλῳ, Κρήτης ἐν πίονι δήμῳ ; Carm. conv.~3 (PMG~885) ap. Ath.~15.694c πλούτου μητέρ᾽ Ὀλυμπίαν ἀείδω / Δήμητρα στεφανηφόροις ἐν ὥραις, / σέ τε, παῖ Διός, Φερσεφόνη· / χαίρετον, εὖ δὲ τάνδ᾽ ἀμφέπετον πόλιν ; \textit{h.Cer.}~486–489 μέγ᾽ ὄλβιος, ὅν τιν᾽ ἐκεῖναι / προφρονέως φίλωνται ἐπιχθονίων ἀνθρώπων· / αἶψα δέ οἱ πέμπουσιν ἐφέστιον ἐς μέγα δῶμα / Πλοῦτον, ὃς ἀνθρώποις ἄφενος θνητοῖσι δίδωσιν ; Pl.~\textit{Cra.}~403a τὸ δὲ Πλούτωνος, τοῦτο μὲν κατὰ τὴν τοῦ πλούτου δόσιν, ὅτι ἐκ τῆς γῆς κάτωθεν ἀνίεται ὁ πλοῦτος, ἐπωνομάσθη ; Σ Ar.~\textit{Plut.}~727 (Dübner, p.~364,~30–41) Τὸν Πλοῦτον Πλούτωνα εἶπε παίζων· ἢ ὅτι καὶ Πλούτωνα αὐτὸν ὑποκοριστικῶς ἐκάλεσεν … τὸν Πλοῦτον οὕτω λέγουσι· καὶ εἰκότως τὸν αὐτὸν τῷ Πλούτωνι τὸν Ἅδην νομίζουσι ; Σ Ar.~\textit{Thesm.}~298 (Dübner, p.~266,~27) ὡς ἀνδρὶ τῆς Περσεφόνης ; Hsch.~ε~7077 εὔπλουτον κανοῦν· εὖ ἔχον πλούτου, διὰ τὰς ἐπ᾽ αὐτῷ ὀλάς (trag. ad. fr. 213 K.-Sn.). πλοῦτον γὰρ ἔλεγον τὴν ἐκ τῶν κριθῶν καὶ τῶν πυρῶν περιουσίαν ; Hsch.~π~2625 πλοῦτος· ἡ ἐκ τῶν σπερμάτων ἐπικαρπία, καὶ ἡ πανσπερμία (ω 486) ; Ar.~\textit{Pax}~1320–1321 κἀπευξαμένους τοῖσι θεοῖσιν / διδόναι πλοῦτον τοῖς Ἕλλησιν ; Soph.~fr. (Inachus) ap. Σ Ar. Plut.~727 Πλούτωνος δ᾽ ἐπείσοδος· … τοιόνδ᾽ ἐμὸν Πλούτων᾽ ἀμεμφίας χάριν ; \textit{Od.}~5.125–127 ὣς δ᾽ ὁπότ᾽ Ἰασίωνι ἐυπλόκαμος Δημήτηρ, / ᾧ θυμῷ εἴξασα, μίγη φιλότητι καὶ εὐνῇ / νειῷ ἔνι τριπόλῳ ; Luc.~\textit{Tim.}~21 ὁ Πλούτων ἀποστέλλει με παρ᾽ αὐτοὺς ἅτε πλουτοδότης καὶ μεγαλόδωρος καὶ αὐτὸς ὤν· δηλοῖ γοῦν καὶ τῷ ὀνόματι ; Hsch.~ε~6720 Εὐβουλεύς· ὁ Πλούτων παρὰ [δὲ] τοῖς πολλοῖς, ὁ ⟨δὲ⟩ Ζεὺς ἐν Κυρήνῃ ; Hsch.~π~2626 Πλούτων· ὁ καταχθόνιος δαίμων (Ar. Ran. 163. 765. 784. Aristophont. fr. 12, 4 K.–A. Plat. Gorg. 523 a) ; Aesch.~\textit{PV}~805–806 οἳ χρυσόρρυτον / οἰκοῦσιν ἀμφὶ νᾶμα Πλούτωνος πόρου ; App. Prov.~4.20 Ὄμπνιος καρπός: ὁ σῖτος καὶ οἱ Δημητριακοὶ καρποί. Ὀμπνία γὰρ ἡ Δήμητρα}
\lp{299 καὶ τῷ Πλούτῳ (τυφλὸς ὁ Πλοῦτος)]}{Macar.~8.60 (CPG II, p.~223) Τυφλὸς ὁ πλοῦτος: ἐπὶ τῶν ἀναξίως εὐπραγούντων ; Apostol.~1.53 (CPG II, p.~253) Ἄζηλος πλοῦτος ὅτι τυφλός: ἐπὶ τῶν κεκτημένων τι καλὸν ἄδηλον ; Tosi n°~106 Fortuna caeca est ; Amphis fr.~23~Kock (\textit{Κουρίς}), ap. Ath.~13.567f τυφλὸς ὁ Πλοῦτος εἶναί μοι δοκεῖ ; Erasm.~\textit{Adag.}~1.7.84 (n°~684) His adscribendum quod Graeci τυφλὸν πλοῦτον appellant, i. e. caecas divitias, reconditas opes \& abstrusas. … Hoc dictum Plutarchus in vita Lycurgi testatur vulgo apud Lacedaemonios fuisse jactatum, quod illic Plutus servaretur inhonoratus, non aliter quam imago picta inanimis \& immobilis}
\lp{299 καὶ τῇ Καλλιγενείᾳ]}{Ar.~fr.~331~K.-A. (Καλλιγένεια) δαίμων περὶ τὴν Δήμητραν, ἣν προλογίζουσαν ἐν ταῖς ἑτέραις Θεσμοφοριαζούσαις ἐποίησεν ; Nonn.~\textit{D.}~6.140 κεῖθι δὲ Καλλιγένειαν, ἑὴν εὔπαιδα τιθήνην ; Plu.~\textit{Quaest.\,Graec.}~31 (298b–c) διὰ τί τοῖς θεσμοφορίοις αἱ τῶν Ἐρετριέων γυναῖκες οὐ πρὸς πῦρ ἀλλὰ πρὸς ἥλιον ὀπτῶσι τὰ κρέα, καὶ Καλλιγένειαν οὐ καλοῦσιν ; Alciphr.~2.37 ἡ μὲν οὖν Ἄνοδος κατὰ τὴν πρώτην γέγονεν ἡμέραν, ἡ Νηστεία δὲ τὸ τήμερον εἶναι παρ᾽ Ἀθηναίοις ἑορτάζεται, τῇ Καλλιγενείᾳ δὲ εἰς τὴν ἐπιοῦσαν θύουσιν ; Σ Ar.~\textit{Thesm.}~80 (Dübner, p.~265,~6–11) ἑνδεκάτη Πυανεψιῶνος ἄνοδος. δωδεκάτη νηστεία, ἐν ᾗ σχολὴν ἄγουσαι ὑπόκεινται αἱ γυναῖκες ἐκκλησιάζουσαι περὶ Εὐριπίδου. τρισκαιδεκάτη καλλιγένεια ; Hsch.~κ~472 Καλλιγένειαν· οὐ τὴν γῆν, ἀλλὰ τὴν Δήμητραν. οὐδεὶς γὰρ οὕτως ἔφη τὴν γῆν καλλιγένειαν. οἱ μὲν τροφὸν αὐτῆς· οἱ δὲ ἱέρειαν· οἱ δὲ ἀκόλουθον (Ar. Thesm. 299) ; Phot.~κ~118 Καλλιγένειαν· Ἀπολλόδωρος μὲν τὴν γῆν· οἱ δὲ Διὸς καὶ Δήμητρος θυγατέρα· Ἀριστοφάνης δὲ ὁ κωμικὸς τροφόν ; Ar.~\textit{Thesm.}~80 ἐπεὶ τρίτη ᾽στὶ Θεσμοφορίων ἡ μέση}
